\documentclass[12pt]{article}
\usepackage[utf8]{inputenc}
\usepackage[T1]{fontenc}
\usepackage{hyperref}
\title{North Mesopotamian Arabic Language Resources for K-12 Curriculum Development}
\author{Compiled by Claude Sonnet 5.5 (claude-sonnet-5-5) for Andr\'as Kornai}
\date{}
\begin{document}
\maketitle

\section{Introduction}
North Mesopotamian Arabic (ISO 639-3 \texttt{ayp}; also called Moslawi, Mardelli or Qeltu Arabic) is a variety of Mesopotamian Arabic spoken in northern Iraq, eastern Syria and south-eastern Turkey. This report is short because the searches found only general and scholarly material, and no school materials.

\section{Primary Educational Resources}
No primer, textbook or school programme specific to this variety was found in the searches. Arabic schooling in the region uses Modern Standard Arabic.

\section{Linguistic and Reference Materials}
\begin{itemize}
\item General overview: North Mesopotamian Arabic, Wikipedia. \url{https://en.wikipedia.org/wiki/North_Mesopotamian_Arabic}
\item Scholarly article on the origin, phonology and loanwords of Moslawi Arabic (Al-Noor Journal for Humanities, 2024). \url{https://jnfh.alnoor.edu.iq/jnfh/article/view/190}
\end{itemize}
No dedicated dictionary or grammar of the dialect turned up in the searches.

\section{Religious and Community Resources}
None found in the searches.

\section{Digital and Audio Resources}
None found in the searches.

\section{Curriculum Development Context}
Little usable material for a curriculum was found. Any curriculum work would need new documentation of the dialect, starting from the scholarly literature cited above, which also lists further works (for example by Aljomard and by Behnstedt and Woidich) that were not checked here.

\end{document}